\chapter{Evaluation}\label{chap:evaluation}
% Budget: 6 standard pages, mostly tables. Every number is from one run of reproduce.sh on 2 November 2026 with
% THREAD_ESCAPE=off (.notes/technical/results-2026-11-02/, with the readings in .notes/author/evaluation-without-concurrency.md),
% except those of Section "Larger Programs", which are from the run of the stages javac and frameworks on 3 November 2026 with
% THREAD_ESCAPE=off (.notes/technical/results-2026-11-03/javac-frameworks/, F88), and where a table says otherwise.

\newcommand{\pct}[1]{\ensuremath{#1\,\%}}

This chapter measures what Graal gains from the facts of our analysis when it compiles a program with GraalVM Native Image,
and what the facts cost. It counts the loads Graal leaves, times the programs, follows the control flow, and measures the
build (\Cref{sec:eval-gain,sec:eval-cost}); then it asks why the gain is as small as it is (\Cref{sec:eval-blockers}) and
whether the images are correct (\Cref{sec:eval-soundness}). The facts are about memory only: no model of concurrency is used,
so every lock and \texttt{volatile} access is assumed to order threads. The commands that regenerate the tables are in
\Cref{app:reproduction}.

\section{Method}\label{sec:eval-method}

An \emph{image} is a native image built from one program by one builder: Native Image from our fork of Graal, extended
with the feature of \Cref{chap:implementation}. The feature takes the facts of Native Image's points-to analysis out of the
build, our tool turns them into kill sets and partitions, and Graal's own phases use them. A \emph{kill set} is the set of
locations a call may write; without one, Graal assumes that the call may write anything and reloads every field after it. A
\emph{configuration} says which facts reach Graal (\Cref{tab:eval-configurations}), and each is compared with the baseline
that differs from it in exactly the facts under test.

\begin{table}[t]
    \centering
    \caption{The configurations. \texttt{graal} hands nothing to Graal, so apart from counting it is Native Image as
    released. \texttt{ceiling} is unsound: an upper bound, not a result.}
    \label{tab:eval-configurations}
    \footnotesize
    \begin{tabular}{llll}
        \toprule
        Configuration & Analysis & Facts given to Graal & Compared with \\
        \midrule
        \texttt{graal}     & default          & none                                   & (baseline) \\
        \texttt{fields}    & default          & kill sets, partitions by field         & \texttt{graal} \\
        \texttt{ceiling}   & default          & kill sets, unsound upper bound         & \texttt{graal}, \texttt{fields} \\
        \texttt{graal}     & allocation sites & none                                   & (baseline) \\
        \texttt{objects}   & allocation sites & kill sets and object-field partitions  & \texttt{graal} on allocation sites \\
        \bottomrule
    \end{tabular}
\end{table}

Native Image runs its points-to analysis in two ways. The \emph{default analysis} names an object by its type; the
\emph{allocation-site analysis}, selected with the builder's option \texttt{AnalysisContextSensitivity=allocsens}, names it by the
place that allocated it. Partitions by object need the finer names; kill sets do not, so we measure them on both. The
allocation-site analysis also changes the code Graal produces with none of our facts: on Are We Fast Yet~\cite{marr2016awfy}
(AWFY) at \texttt{-O2} the calls left fall from 1,407 to 1,110 and the loads rise from 943 to 954. We therefore compare images
only on the same analysis, and report the price of the allocation-site analysis separately (\Cref{sec:eval-cost}).

\begin{table}[t]
    \centering
    \caption{The programs.}
    \label{tab:eval-programs}
    \footnotesize
    \begin{tabular}{p{3.0cm}p{6.8cm}p{2.5cm}}
        \toprule
        Program & What it is & Used for \\
        \midrule
        \texttt{LoadStoreElimination}   & 34 methods: a call that writes little between a store and a load; 12 of them are controls that must keep their loads & loads \\
        \texttt{Signatures} & 17 methods, one per rule for what a call's signature says it may touch & loads \\
        \texttt{Aliasing}   & 13 methods: two references that may or may not meet, with 4 controls & loads \\
        \texttt{Allocation} & 4 methods: a call that allocates enough to collect, which tests the virtual machine state a call kills & loads \\
        AWFY                & 14 benchmarks over a small collection library of their own, 459 methods, pinned at commit \texttt{74306fe}; a \texttt{main} constructs each benchmark and prints whether it verified & loads, run time, cost \\
        javac               & the JDK's compiler, compiling the 79 sources of AWFY & cost, loads, run time \\
        PetClinic           & Graal's Native Image bundle of Spring PetClinic (Tomcat, H2, HikariCP), served under load & cost, loads, run time \\
        Barista             & eleven server applications of the Barista suite, release 0.6.6: a plain hello world, hello worlds on seven frameworks, two more Micronaut applications, and Spring PetClinic & cost, loads \\
        Native Image's tests & its own JUnit suite: 23 images with 289 tests & robustness \\
        \bottomrule
    \end{tabular}
\end{table}

The four small programs were written to contain calls that write little, so the answer is known: they show that the
mechanism works, not what a typical program gains. AWFY is the realistic program.

\begin{table}[t]
    \centering
    \caption{Conditions of the measurements.}
    \label{tab:eval-conditions}
    \footnotesize
    \begin{tabular}{p{2.6cm}p{9.6cm}}
        \toprule
        Item & Setting \\
        \midrule
        Machine   & the faculty server pontos04: two AMD EPYC 7371, 64 threads, 1\,TB of RAM, Debian 13; the only machine, without hardware counters (\texttt{perf\_event\_paranoid} is 3) \\
        Builder   & GraalVM 25.5-dev, built from our fork of Graal on a JDK~25 with JVMCI \\
        Run       & 2 November 2026, one run of \texttt{reproduce.sh} with \texttt{THREAD\_ESCAPE=off}; tool \texttt{df40ed9}, fork \texttt{12284dadc70} (control flow: \texttt{acb8880}, \texttt{8a70f22a6ac}) \\
        AWFY builds & \texttt{-O2} as the analysis builds it, and \texttt{-O3} with trivial inlining as a user would; they are different compilations, so only the columns of one build compare \\
        Correctness & every image must print exactly what the JVM prints; a build whose output differs is a failure, not a data point \\
        Counts    & field loads, stores, and calls left in the program's methods at the end of Graal's low tier; static counts, which loop unrolling multiplies, not run time \\
        Run time  & 14 benchmarks, 20 iterations, five runs, the first iteration dropped, the median; four cores of one NUMA node, nothing else running; images take turns run by run \\
        Builds timed & \texttt{graal} and \texttt{fields} four times each, timed together; the first \texttt{graal} binary timed twice (A/A) shows the noise \\
        Larger programs & 3 November 2026, the stages \texttt{javac} and \texttt{frameworks} with \texttt{THREAD\_ESCAPE=off}; tool \texttt{f185b31}, fork \texttt{308e7320996}; three builds of \texttt{graal} and of \texttt{fields} of each, one at a time \\
        Build time & wall-clock seconds of one build at a time on an empty machine, configurations taking turns, five rounds; median and range, stage times from the builder's log \\
        \bottomrule
    \end{tabular}
\end{table}

\section{What Graal Optimizes with the Facts}\label{sec:eval-gain}

\subsection{Loads Left}\label{sec:eval-loads}

\begin{table}[t]
    \centering
    \caption{Loads left in the program's methods; the percentage is the change against \texttt{graal} on the same analysis.}
    \label{tab:eval-loads}
    \footnotesize
    \begin{tabular}{lrrrrr}
        \toprule
        & \multicolumn{2}{c}{Default analysis} & \multicolumn{3}{c}{Allocation sites} \\
        \cmidrule(lr){2-3} \cmidrule(lr){4-6}
        Program & \texttt{graal} & \texttt{fields} & \texttt{graal} & \texttt{fields} & \texttt{objects} \\
        \midrule
        \texttt{LoadStoreElimination}   & 40 & 27 (\pct{-32.5}) & 43 & 30 (\pct{-30.2}) & 30 (\pct{-30.2}) \\
        \texttt{Signatures} & 16 & 12 (\pct{-25.0}) & 16 & 12 (\pct{-25.0}) & 12 (\pct{-25.0}) \\
        \texttt{Aliasing}   & 16 & 14 (\pct{-12.5}) & 18 & 16 (\pct{-11.1}) & 13 (\pct{-27.8}) \\
        \texttt{Allocation} &  9 &  7 (\pct{-22.2}) &  9 &  7 (\pct{-22.2}) &  7 (\pct{-22.2}) \\
        \midrule
        Four small programs & 81 & 60 (\pct{-25.9}) & 86 & 65 (\pct{-24.4}) & 62 (\pct{-27.9}) \\
        \midrule
        AWFY, \texttt{-O2}  & 943 & 903 (\pct{-4.2}) & 954 & 914 (\pct{-4.2}) & 914 (\pct{-4.2}) \\
        AWFY, \texttt{-O3}  & 2,565 & 2,511 (\pct{-2.1}) & -- & -- & -- \\
        \bottomrule
    \end{tabular}
\end{table}

\begin{table}[t]
    \centering
    \caption{The default-analysis loads of the small programs, \texttt{graal} to \texttt{fields}, by kind of method.
    \emph{Controls} are the methods named \texttt{...Control}: a call in them may write the field, or a reference may alias
    another, so they must keep their loads. \emph{Unchanged} are the loads of the other methods that do not change (helpers,
    \texttt{main}, inlined monitor code, and in \texttt{Signatures} the three methods that cross a monitor, a native call, and a
    \texttt{volatile} access). \emph{Reduced} are the loads of the methods that lose loads. The three add up to the totals of
    the default analysis in the table above.}
    \label{tab:eval-loads-split}
    \footnotesize
    \begin{tabular}{lrrr}
        \toprule
        Program & Controls & Unchanged & Reduced \\
        \midrule
        \texttt{LoadStoreElimination}   & $16 \to 14$ & 10 & $14 \to 3$ \\
        \texttt{Signatures} & --          & 12 & $4 \to 0$ \\
        \texttt{Aliasing}   & $5 \to 4$   &  9 & $2 \to 1$ \\
        \texttt{Allocation} & --          &  0 & $9 \to 7$ \\
        \midrule
        Four small programs & $21 \to 18$ & 31 & $29 \to 11$ (\pct{-62.1}) \\
        \bottomrule
    \end{tabular}
\end{table}

On the four small programs, kill sets remove 12.5--32.5\,\% of the loads, and \pct{62} (29 to 11) in the methods that lose
loads at all; the controls keep theirs, as they must. On AWFY they remove 40 of 943 loads (\pct{4.2}) at \texttt{-O2} and 54
of 2,565 (\pct{2.1}) at \texttt{-O3}, and the same 40 on allocation sites. Partitions by allocation site add nothing on AWFY:
\texttt{objects} leaves exactly the loads that \texttt{fields} leaves, because they split one of its 234 fields. They pay where
the program was written for them: on \texttt{Aliasing} they remove three loads more (16 to 13).

\subsection{Run Time}\label{sec:eval-time}

A load that is no longer reloaded saves little by itself; it pays where it enables something more. We therefore time the
programs. \Cref{tab:eval-time} gives AWFY at \texttt{-O3} on the default analysis.

\begin{table}[t]
    \centering
    \caption{Run time of AWFY at \texttt{-O3}, default analysis, in percent against the first \texttt{graal} build (negative is faster).}
    \label{tab:eval-time}
    \footnotesize
    \begin{tabular}{lrrrr}
        \toprule
        Benchmark & \texttt{graal}, ms & A/A & other \texttt{graal} builds & \texttt{fields}, four builds \\
        \midrule
        Queens      & 23.3  & $+0.4$ & $-1.8$ to $-1.5$ & $\mathbf{-27.8}$ \textbf{to} $\mathbf{-27.5}$ \\
        Sieve       & 84.9  & $-0.8$ & $-5.3$ to $-4.3$ & $-11.3$ to $-8.1$ \\
        Storage     & 42.3  & $-0.4$ & $-0.2$ to $+0.3$ & $-4.8$ to $-2.7$ \\
        Permute     & 27.1  & $+1.1$ & $-1.1$ to $+0.4$ & $+7.2$ to $+8.5$ \\
        List        & 39.5  & $0.0$  & $+1.3$ to $+2.3$ & $+1.5$ to $+4.1$ \\
        Richards    & 97.6  & $+0.6$ & $+3.3$ to $+9.9$ & $+1.7$ to $+11.5$ \\
        DeltaBlue   & 14.4  & $-0.3$ & $+0.2$ to $+0.9$ & $-0.1$ to $+1.7$ \\
        Havlak      & 65.7  & $+0.6$ & $+1.5$ to $+2.5$ & $-2.1$ to $+2.4$ \\
        Other six   & 27--83 & $-0.1$ to $+0.3$ & $-0.4$ to $+1.3$ & $-0.5$ to $+1.2$ \\
        \midrule
        Total       & 699.6 & $+0.1$ & $0.0$ to $+1.1$ & $-1.4$ to $-0.1$ \\
        \bottomrule
    \end{tabular}
\end{table}

The six benchmarks not listed are Bounce, CD, Json, Towers, Mandelbrot, and NBody; Mandelbrot and NBody do not change.
\textbf{Queens} is the one clear win: \pct{-27.5} to \pct{-27.8} in every build, far outside the noise. Its recursive method
\texttt{placeQueen} no longer has each call kill four array-reference fields, so their loads leave its loop
(\Cref{sec:eval-control}). \textbf{The total} is \pct{-0.1} to \pct{-1.4} (mean \pct{-0.8}), against \pct{0.0} to \pct{+1.1} for
the other builds of \texttt{graal}: about 1\,\%, mostly Queens. Sieve and Storage are faster in all four builds, and we did not
investigate why. \textbf{Permute} is \pct{+7.2} to \pct{+8.5} slower in every build. Its code changes (\Cref{sec:eval-control}), but
in an earlier experiment that moved all methods to a 16, 64, or 128-byte alignment instead of 32, \texttt{fields} was 6.5--7.1\,\%
\emph{faster} than \texttt{graal}, so we read it as code placement and not as an effect of the facts. Richards, List, Havlak, and
DeltaBlue move by several percent between builds of one configuration, and their ranges for \texttt{fields} overlap those of the
other builds of \texttt{graal}; we make no claim about them.

\subsection{The Ceiling}\label{sec:eval-ceiling}

To learn whether more precise facts would help, we built the \texttt{ceiling} image, which leaves out the kill of everything the
facts cannot name: identity hash codes, class initialization, synchronization, native and unresolved calls. It is unsound by
construction, but it prints what the JVM prints on AWFY, and it is the most that precise kill sets could give Graal.

\begin{table}[t]
    \centering
    \caption{The ceiling on AWFY at \texttt{-O3}, default analysis: the loads left and the run time in percent against
    \texttt{graal}, one build of each; the loads of the four builds timed together differ from these by at most one.}
    \label{tab:eval-ceiling}
    \footnotesize
    \begin{tabular}{lrrrrrr}
        \toprule
        & Loads & In loops & Total & Queens & Permute & Richards \\
        \midrule
        \texttt{graal}   & 2,565 & 937 & --      & --       & --      & -- \\
        \texttt{fields}  & 2,511 (\pct{-2.1}) & 904 & $+0.4$ & $-27.9$ & $+6.5$ & $+13.7$ \\
        \texttt{ceiling} & 2,393 (\pct{-6.7}) & 859 & $-1.5$ & $-26.7$ & $-7.6$ & $-0.7$ \\
        \bottomrule
    \end{tabular}
\end{table}

The ceiling removes three times as many loads as \texttt{fields} (172 against 54) and is not clearly faster: Queens, the one
benchmark that gains, gains the same with \texttt{fields}, and Permute and Richards differ by several percent between builds of
one configuration anyway. Making the remaining opaque calls precise would not make AWFY faster.

\subsection{Control Flow and Loops}\label{sec:eval-control}

A load that no longer depends on a call may leave a loop, and a branch that depended on it may fold. We count the blocks,
control-flow edges, ifs, and loops of every method at the end of Graal's low tier, and for every load whether it is in a loop.
The images are built with source positions, so that the same load can be found in two images.

\begin{table}[t]
    \centering
    \caption{Control flow of AWFY at \texttt{-O3}, default analysis, summed over the 185 methods the images compile.}
    \label{tab:eval-control}
    \footnotesize
    \begin{tabular}{lrr}
        \toprule
        & \texttt{graal} & \texttt{fields} \\
        \midrule
        Blocks                       & 21,197 & 21,201 \\
        Control-flow edges           & 28,997 & 29,003 \\
        Ifs                          &  8,515 &  8,514 \\
        Loops                        &    602 &    599 \\
        Methods with other control flow & -- & 9 \\
        \bottomrule
    \end{tabular}
\end{table}

\begin{table}[t]
    \centering
    \caption{The methods whose control flow changes most, \texttt{graal} to \texttt{fields}. The other five of the nine change
    by at most 12 edges each.}
    \label{tab:eval-control-methods}
    \footnotesize
    \begin{tabular}{lrrrr}
        \toprule
        Method & Blocks & Edges & Ifs & Loops \\
        \midrule
        \texttt{Queens.placeQueen}  & $85 \to 207$ & $119 \to 299$ & $39 \to 97$ & $1 \to 0$ \\
        \texttt{Permute.benchmark}  & $106 \to 15$ & $152 \to 19$  & $49 \to 5$  & $1 \to 0$ \\
        \texttt{Permute.permute}    & $31 \to 16$  & $41 \to 19$   & $13 \to 6$  & -- \\
        \texttt{Queens.queens}      & $120 \to 126$ & $169 \to 176$ & --          & $2 \to 1$ \\
        \bottomrule
    \end{tabular}
\end{table}

\begin{table}[t]
    \centering
    \caption{The loads that \texttt{fields} removes from AWFY at \texttt{-O3}, by where they were. A load that left a loop and
    appears outside it was hoisted, peeled, or unrolled: the counts cannot tell which. The numbers are changes in the count of
    loads, \texttt{graal} to \texttt{fields}.}
    \label{tab:eval-loops}
    \footnotesize
    \begin{tabular}{lr}
        \toprule
        Loads & Change \\
        \midrule
        in loops                                               & $-33$ \\
        \quad of these, left the loop and appear outside it    & 5 \\
        \quad of these, gone (forwarded, merged, or folded)    & 28 \\
        \quad entered a loop                                   & 0 \\
        outside loops (the 5 included)                         & $-21$ \\
        \midrule
        all loads                                              & $-54$ \\
        \bottomrule
    \end{tabular}
\end{table}

Kill sets barely change the control flow as a whole (6 edges of 29,000), but they change it a lot in three methods. In
\texttt{Queens.placeQueen} the loop, which contains the recursive call, loses its six loads of the four array-reference fields
and then disappears; the method is 2.4 times larger (85 to 207 blocks) and has more edges, because its ifs depend on data. In
\texttt{Permute.benchmark} the loop disappears too, and most branches with it (49 to 5 ifs); we did not establish why. We read the
loops as unrolled by Graal once the loads were gone, but we did not read its decision. Loop-invariant code motion in the strict
sense is rare: of the 33 loads that left a loop, 5 appear outside it and 28 are gone.

\section{Cost}\label{sec:eval-cost}

\subsection{AWFY}\label{sec:eval-cost-awfy}

\begin{table}[t]
    \centering
    \caption{Build of AWFY at \texttt{-O3}, median of five rounds in seconds (range of the total).}
    \label{tab:eval-cost}
    \footnotesize
    \setlength{\tabcolsep}{4pt}
    \begin{tabular}{lrrrr}
        \toprule
        Configuration & Total & Analysis & Compile & Memory, GiB \\
        \midrule
        \texttt{graal}                          & 24.1 (23.1--24.8) &  4.4 & 7.4 & 5.3 \\
        \texttt{fields}                         & 25.4 (24.8--25.6) &  6.3 & 7.7 & 5.4 \\
        \quad added                             & $+1.3$ (\pct{+5}) & $+1.9$ (\pct{+43}) & $+0.3$ (\pct{+4}) & $+0.1$ \\
        \midrule
        \texttt{graal} on allocation sites      & 30.4 (30.3--32.6) & 11.8 & 7.4 & 5.5 \\
        \texttt{objects} on allocation sites    & 33.3 (32.6--34.1) & 13.7 & 7.8 & 5.5 \\
        \quad added                             & $+2.9$ (\pct{+10}) & $+1.9$ (\pct{+16}) & $+0.4$ (\pct{+5}) & 0.0 \\
        \bottomrule
    \end{tabular}
\end{table}

Kill sets add 1.3 seconds (\pct{5}) to the build of AWFY, 1.9 of them in the analysis stage, where the facts are exported and
the tool runs. The compile stage does not move (0.3 seconds against a spread of 6.9--8.1), and neither does the memory. The
allocation-site analysis costs more than everything we add: \texttt{graal} on allocation sites takes 30.4 seconds against 24.1
(\pct{+26}) and its analysis stage 2.7 times as long, and the partitions then add 2.9 seconds (\pct{+10}) and gave nothing on
AWFY. The price of the finer analysis grows with the program: javac builds in 67 seconds on the default analysis and in 288 on
allocation sites (24 October, three rounds, nothing of ours in either).

\subsection{Larger Programs}\label{sec:eval-large}

AWFY has 459 methods; the programs for which Native Image is built have many more. \Cref{tab:eval-large} gives the build of
javac, of the PetClinic bundle, and of the eleven Barista applications on the default analysis, \texttt{graal} against
\texttt{fields}, three builds of each, one at a time on an empty machine, the configurations taking turns.%
\footnote{The Barista suite is at \url{https://github.com/barista-benchmarks/barista}. Its two Quarkus applications are left out:
they do not build with Native Image from our fork even without our feature.} In every program the ranges of the two
configurations do not overlap, so three builds tell the increment from the spread.

\begin{table}[t]
    \centering
    \caption{Build of the larger programs, median of three builds in seconds (range of the totals), by the build of \texttt{graal}.}
    \label{tab:eval-large}
    \footnotesize
    \setlength{\tabcolsep}{3pt}
    \begin{tabular}{lrrrrr}
        \toprule
        Program & \texttt{graal} & \texttt{fields} & Added & Analysis & Loads \\
        \midrule
        Spring PetClinic & 128.0 (127.0--130.0) & 172.0 (170.0--173.0) & +44.0 (\pct{+34}) & +42.6 & \pct{-0.9} \\
        PetClinic & 116.0 (112.0--119.0) & 146.0 (146.0--151.0) & +30.0 (\pct{+26}) & +33.8 & \pct{-0.9} \\
        play-scala & 67.0 (67.0--68.0) & 82.0 (79.0--84.0) & +15.0 (\pct{+22}) & +14.9 & \pct{-0.8} \\
        dropwizard & 63.0 (62.0--63.0) & 73.0 (72.0--80.0) & +10.0 (\pct{+16}) & +12.2 & \pct{-0.9} \\
        javac & 62.0 (61.0--64.0) & 76.0 (69.0--78.0) & +14.0 (\pct{+23}) & +7.9 & \pct{-3.1} \\
        spring-hello & 62.0 (60.0--64.0) & 73.0 (71.0--74.0) & +11.0 (\pct{+18}) & +11.4 & \pct{-0.8} \\
        micronaut-shopcart & 61.0 (60.0--62.0) & 72.0 (71.0--75.0) & +11.0 (\pct{+18}) & +10.4 & \pct{-0.6} \\
        micronaut-similarity & 56.5 (56.2--57.3) & 66.0 (64.0--67.0) & +9.5 (\pct{+17}) & +10.0 & \pct{-0.7} \\
        ktor & 55.8 (55.2--57.6) & 65.0 (64.0--67.0) & +9.2 (\pct{+16}) & +10.1 & \pct{-0.4} \\
        micronaut-hello & 55.5 (55.5--56.3) & 66.0 (65.0--67.0) & +10.5 (\pct{+19}) & +9.6 & \pct{-0.7} \\
        helidon & 46.0 (44.9--46.5) & 52.8 (51.0--52.8) & +6.8 (\pct{+15}) & +5.9 & \pct{-0.8} \\
        vertx & 37.3 (36.0--37.9) & 41.1 (40.6--43.5) & +3.8 (\pct{+10}) & +5.1 & \pct{-0.8} \\
        plain hello world & 29.9 (29.1--30.2) & 32.1 (31.7--32.5) & +2.2 (\pct{+7}) & +2.6 & -- \\
        \bottomrule
    \end{tabular}
\end{table}

The plain hello world has 3 loads in its own code. Kill sets add 7 to 34\,\% to the build: 2.2 seconds to the plain hello world, 6.8 to 15.0 seconds (\pct{+15} to \pct{+22}) to
the eight framework applications of ordinary size, and 30 and 44 seconds (\pct{+26} and \pct{+34}) to the two PetClinics, the
bundle of Graal and the one of Barista. In every program the percentage is larger than on AWFY (\pct{+5}), and it grows with the
program. The increment
is in the analysis stage, where the facts are exported and the tool runs; the compile stage moves by \mbox{$-1.9$} to
$+0.4$ seconds. The analysis column can exceed the total: on PetClinic the stage grows by 33.8 seconds and the build by 30.0,
because the other stages together are 3.8 seconds shorter in the \texttt{fields} builds. That is within their spread between
builds of one configuration (the compile stage of the three \texttt{graal} builds took 20.9--23.1 seconds, of the three
\texttt{fields} builds 20.4--21.7). The peak memory of the servers differs by at most 0.3 GiB; javac's rises from 6.2 to 7.6
GiB, AWFY's by 0.1.

\textbf{javac} is the exception in how the increment divides. Its analysis stage grows by 7.9 seconds (\pct{+13} of the
build), but the total by 14.0: the rest is in the stage that lays out the methods, which took 10.4--10.7 seconds in every
\texttt{graal} build and in one of the \texttt{fields} builds, and 16.0 and 16.4 seconds in the other two. Three builds do not
tell whether \texttt{fields} slows that stage or some builds are slow; the build with the fast stage took 69 seconds against
61--64 for \texttt{graal}, which is 7 seconds more. The analysis increment of 7.9 seconds is the one to cite.

The loads that \texttt{fields} removes are few: 0.4 to 0.9\,\% of the loads in the servers' own methods, and 3.1\,\% in javac
(7,846 to 7,531 of them in loops). The calls left do not change (javac 38,299 to 38,301, PetClinic 278,122 to 278,171). We did
not establish why the servers gain so little.

\textbf{Run time} can be told for two of the programs. javac compiles the 79 sources of AWFY in 0.240--0.241 seconds in the
three \texttt{graal} builds and in 0.239--0.243 in the three \texttt{fields} builds (ten runs of five compilations each); the
first \texttt{graal} image timed again differs by $-0.3$\,\%, so there is no difference. PetClinic served 4,085--4,344 requests
per second in the three \texttt{graal} builds and 4,175--4,384 in the three \texttt{fields} builds (16 connections, five runs
of 30 seconds each after a warm-up of 15); the first \texttt{graal} image measured twice differs by \pct{+0.4}, the runs of one
image by about \pct{1}, but the builds of one configuration by \pct{6.3} (\texttt{graal}) and \pct{5.0} (\texttt{fields}). The
medians differ by \pct{+2.2}, which is within the spread between builds, so it is not a gain. No Barista application was timed:
with at most \pct{0.9} fewer loads there is nothing to see through the spread between builds.

\subsection{Allocation Sites}\label{sec:eval-cost-sites}

Partitions by object need Native Image's allocation-site analysis, which costs Native Image itself. \Cref{tab:eval-sites} gives
that price from the builds we have: earlier runs, each row from one run, with the code and the model of concurrency of its day,
so the numbers do not compare with the other tables. The allocation sites multiply the build by 1.3 on AWFY and by 4.3 on javac.
On PetClinic a build with allocation sites and a cutoff of 100, with nothing of ours in it, took 1,830 seconds against 120, 15
times as long. On the ten Barista applications that built, an image with allocation sites, the cutoff, and our export and lock
elision took 3.0 to 10.7 times the \texttt{graal} build; where we timed our steps (three applications), they were 6 to 17
seconds of it, so the factor is the price of the allocation sites. Spring PetClinic did not finish: our export ran out of
memory. The factor grows with the program, and it is far above what the facts themselves cost (7 to 34\,\%).

\begin{table}[t]
    \centering
    \caption{Build time of Native Image on the default analysis and with allocation sites, in seconds, and the ratio.}
    \label{tab:eval-sites}
    \footnotesize
    \begin{tabular}{lrrr}
        \toprule
        Program (builds) & Default & Allocation sites & Times \\
        \midrule
        AWFY (five rounds)                  & 24.1 & 30.4 & 1.3 \\
        javac (three rounds)               & 67   & 288  & 4.3 \\
        PetClinic (one build)              & 120  & 1,830 & 15 \\
        \midrule
        Spring PetClinic (one build)       & 138 & out of memory & -- \\
        play-scala                          & 72   & 771  & 10.7 \\
        dropwizard                          & 69   & 537  & 7.8 \\
        micronaut-shopcart                  & 66   & 647  & 9.8 \\
        spring-hello                        & 65   & 567  & 8.7 \\
        ktor                                & 63   & 364  & 5.8 \\
        micronaut-hello                     & 62   & 550  & 8.9 \\
        micronaut-similarity                & 59   & 460  & 7.8 \\
        helidon                             & 50   & 377  & 7.5 \\
        vertx                               & 42   & 265  & 6.3 \\
        plain hello world                   & 33   & 99   & 3.0 \\
        \bottomrule
    \end{tabular}
\end{table}

\subsection{Where the Time Goes}\label{sec:eval-cost-steps}

The added time is our work in the analysis stage. \Cref{tab:eval-steps} divides it into its steps: the facts are collected from
Native Image's analysis (together with a text report of them, which these builds discard), the facts are written to the export
file, and the tool computes the kill sets from it. The three add up to the increment of the analysis stage in every program, to within 0.8 seconds. Collecting the facts is the
largest step in all programs but the two PetClinics, where writing the export is (14.9 of 33.6 seconds on PetClinic, 17.6 of
42.5 on Spring PetClinic); the tool takes 15 to 30\,\%. Applying the kill sets happens in compilation, on several threads, and
adds at most 2.3 seconds of thread time. The last row is AWFY built with partitions by object on allocation sites, against
\texttt{graal} on allocation sites: its steps are as large as those of \texttt{fields}, and the check that every graph keeps
its partitions consistent adds 0.4 seconds of thread time. No other program was built with partitions by object in this run,
so no other program has such a row.

\begin{table}[t]
    \centering
    \caption{Our steps in the \texttt{fields} builds, and in the \texttt{objects} build of AWFY, in seconds (median).}
    \label{tab:eval-steps}
    \footnotesize
    \setlength{\tabcolsep}{4pt}
    \begin{tabular}{lrrrrrr}
        \toprule
        & \multicolumn{3}{c}{Steps} & & & \\
        \cmidrule(lr){2-4}
        Program & Facts & Write & Tool & Sum & Stage added & Apply \\
        \midrule
        Spring PetClinic & 12.3 & 17.6 & 12.7 & 42.5 & +42.6 & 2.3 \\
        PetClinic & 9.8 & 14.9 & 8.9 & 33.6 & +33.8 & 1.7 \\
        play-scala & 6.0 & 5.1 & 3.4 & 14.5 & +14.9 & 1.1 \\
        dropwizard & 6.0 & 3.2 & 2.3 & 11.5 & +12.2 & 0.6 \\
        javac & 4.8 & 1.4 & 1.2 & 7.4 & +7.9 & 0.5 \\
        spring-hello & 5.6 & 3.4 & 2.5 & 11.5 & +11.4 & 0.7 \\
        micronaut-shopcart & 4.7 & 3.7 & 2.6 & 11.0 & +10.4 & 0.9 \\
        micronaut-similarity & 4.2 & 2.9 & 2.1 & 9.3 & +10.0 & 0.7 \\
        ktor & 5.0 & 2.5 & 1.9 & 9.4 & +10.1 & 0.6 \\
        micronaut-hello & 4.2 & 2.9 & 2.1 & 9.2 & +9.6 & 0.7 \\
        helidon & 3.1 & 1.8 & 1.2 & 6.0 & +5.9 & 0.3 \\
        vertx & 3.3 & 1.0 & 0.9 & 5.2 & +5.1 & 0.3 \\
        plain hello world & 1.9 & 0.5 & 0.5 & 2.9 & +2.6 & 0.1 \\
        \midrule
        AWFY & 1.5 & 0.2 & 0.3 & 1.9 & +1.9 & 0.1 \\
        AWFY on allocation sites, \texttt{objects} & 1.5 & 0.2 & 0.3 & 2.0 & +1.9 & 0.3 \\
        \bottomrule
    \end{tabular}
\end{table}

\section{What Blocks the Analysis}\label{sec:eval-blockers}

The gain on AWFY is small because most of its methods remain opaque: of the 459 methods of the program, 199 still kill
everything. \Cref{tab:eval-blockers} shows what leaving out the kill of each reason would free.

\begin{table}[t]
    \centering
    \caption{Methods of AWFY (of 459) freed from killing everything when the kill of the named effects is left out of the
    kill sets (export of 22 October).}
    \label{tab:eval-blockers}
    \footnotesize
    \begin{tabular}{p{0.7\linewidth}r}
        \toprule
        Left out & Methods freed \\
        \midrule
        identity hash codes, synchronization, unresolved calls, or other effects, each alone & 0 \\
        class initialization alone & 12 \\
        identity hash codes, class initialization, and synchronization together & 12 \\
        \dots{} and unresolved calls & 12 \\
        \dots{} and every other effect & 199 \\
        \bottomrule
    \end{tabular}
\end{table}

No single reason is enough. Each of the 187 methods that remain after the first three reaches all of an unresolved call, a
transition to native code, a raw write, and a snippet, and the reasons form a chain: a method that may build an exception, or
call \texttt{equals} on an object of unknown type, reaches a lock, then class initialization, then a native call. The layers have
to go together, and even then AWFY does not run faster (\Cref{tab:eval-ceiling}).

\section{Soundness in Practice}\label{sec:eval-soundness}

\Cref{chap:effects} argues soundness for the model; the proofs are sketches in three places. The evidence of this chapter is of
another kind, and \Cref{tab:eval-checks} lists it.

\begin{table}[t]
    \centering
    \caption{Evidence that the images are correct.}
    \label{tab:eval-checks}
    \footnotesize
    \begin{tabular}{p{3.2cm}p{9.0cm}}
        \toprule
        Check & Result \\
        \midrule
        Output                    & every image of this chapter prints exactly what the JVM prints \\
        Controls that must fail   & a method claimed to write nothing gives 171 instead of 177 in \texttt{LoadStoreElimination} at \texttt{-O3}, so a false claim does not pass unseen \\
        Native Image's own tests  & with the feature and \texttt{-J-ea -J-esa}, its 23 JUnit images build and all 289 tests pass; the assertions found three defects that the examples did not: disabled flows asked for their objects, calls killed the initialization state, and a JSON escape failed on half a surrogate pair \\
        \bottomrule
    \end{tabular}
\end{table}

This is evidence by controls and examples, not a proof.

\section{Discussion}\label{sec:eval-discussion}

\paragraph{Benefit.} Graal optimizes a little more with the facts, and in a specific way. Where it must assume that a call writes
anything and the call sits in a loop, the kill set lets loads leave the loop and the loop change shape; Queens is that case,
with \pct{-27.7}. Elsewhere the gain is a handful of loads (2--4\,\% of AWFY's) that shows in the total as about 1\,\%. The
small size is not a weakness of the implementation: the ceiling removes three times the loads and is no faster, and the
opacity of AWFY has several layers, each of which would have to go.

\paragraph{Cost.} The facts are cheap on AWFY: kill sets add \pct{5} to the build. On larger programs they add 7 to 34\,\%
(\Cref{tab:eval-large}), almost all of it in the analysis stage, and they remove 0.4 to 3.1\,\% of the loads. The allocation-site analysis that partitions by object need is dearer still
(\Cref{tab:eval-sites}): it adds \pct{26} to the build of AWFY, multiplies that of javac by 4.3 and that of PetClinic by 15,
and takes the Barista applications 3.0 to 10.7 times as long; the partitions gained nothing on AWFY. Kill sets on the default
analysis are therefore the affordable combination.

\paragraph{Not shown.} We did not time the Barista applications, and PetClinic's throughput differs less between the
configurations than between builds of one, so the larger programs show a cost and no run-time gain. We did not count loads at
run time.

\section{Threats to Validity}\label{sec:eval-threats}

\paragraph{One machine and few programs.} Every timing is from one server with an AMD Zen~1 processor and without hardware
counters. The realistic program for run time is AWFY, with javac and PetClinic as larger checks; the gains on the four small
programs are by design.

\paragraph{Code placement.} Native Image places every method at a multiple of 32 bytes, so a change in the size of one method
moves all that follow it. Placement alone moves twelve of the fourteen AWFY benchmarks by 4--13\,\% with identical instructions,
which is above the effect of the facts for all but Queens. We address it by timing four builds of each configuration together
and by not citing the benchmarks that vary between builds.

\paragraph{Static counts.} A load left is not a load executed, and loop unrolling multiplies counts; \texttt{Allocation}'s 66
stores are one store unrolled 64 times. That Graal unrolled the loops of \texttt{Queens.placeQueen} and \texttt{Permute.benchmark} is
our reading of the counts.

\paragraph{One wrong baseline.} In the first run of the control-flow measurement one \texttt{graal} baseline (AWFY on
allocation sites at \texttt{-O3}) printed an incorrect result on every run of the binary. Four rebuilds of it and a complete
second run, 25 images, were correct. That configuration hands nothing to Graal, so our facts are not the cause; we have not
explained it. The tables use the second run.

\paragraph{Single builds and older numbers.} The ceiling is one build of each configuration, three builds of each configuration
are all that the larger programs have, the allocation-site builds of \Cref{tab:eval-sites} are of 24 to 29 October (one build each for PetClinic and Barista), and the export behind \Cref{sec:eval-blockers} is of 22 October. The \texttt{-O3} build on allocation sites was not run.
